\documentclass{article}
\usepackage[utf8]{inputenc}

\title{Computación Distribuida \\ Tarea 06}
\author{Olea Olvera Marco Iván \\ Arteaga Hernández Alejandro Iabin}
\date{}

\usepackage[top=1.5cm, bottom=1.5cm, left=2cm, right=2cm]{geometry}
\usepackage{fouriernc}
\usepackage{amssymb}
\usepackage{listings}
\lstset{
  basicstyle=\ttfamily,
  mathescape
}

\usepackage{graphicx}
\graphicspath{{./}}

\begin{document}

\maketitle

\section{}

Sean $\phi_1$ y $\phi_2$ dos eventos en $\alpha$.

\begin{itemize}

    \item Supongamos que $\phi_1 \rightarrow \phi_2$. Queremos ver que $VC(\phi_1) < VC(\phi_2)$. Por la definición de la relación $\rightarrow$ (sucedió antes que), tenemos tres casos:
    \begin{enumerate}
        \item $\phi_1$ y $\phi_2$ son eventos del mismo procesador $p_i$ y $\phi_1$ ocurre antes que $\phi_2$ en $\alpha$. \\
        Ninguna entrada de $VC_i$ puede decrecer, es decir, $VC(\phi_1)[j] \leq VC(\phi_2)[j]$ para toda $j$, y esto implica que $VC(\phi_1) \leq VC(\phi_2)$. Además, en particular, $VC(\phi_1)[i] < VC(\phi_2)[i]$ ya que $VC_i[i]$ incrementa en uno con cada evento de $p_i$, y $\phi_1$ ocurre antes que $\phi_2$. Por lo tanto, $VC(\phi_1) \not= VC(\phi_2)$ y esto significa que $VC(\phi_1) < VC(\phi_2)$.
        \item $\phi_1$ es el evento en el que se manda un mensaje $m$ de $p_i$ a $p_j$, y $\phi_2$ es el evento en que $p_j$ recibe $m$. \\
        Durante $\phi_2$, $p_j$ actualiza su vector de la siguiente manera: \\
        $VC(\phi_2)[j] = max\{VC(\phi_1)[j],\ VC_j[j]\}$ \\
        donde $VC_j$ es el vector de $j$ antes del evento $\phi_2$. Es decir, $VC(\phi_1) \leq VC(\phi_2)$. Luego, $p_j$ incrementa la j-ésima entrada del vector en uno, por lo que $VC(\phi_1) \not= VC(\phi_2)$. Por lo tanto, $VC(\phi_1) < VC(\phi_2)$.
        \item Existe un evento $\phi$ tal que $\phi_1 \rightarrow \phi$ y $\phi \rightarrow \phi_2$. \\
        Sin importar si $\phi$ es un evento del mismo procesador que $\phi_1$ (primer caso) o es el evento en el que otro procesador recibió un mensaje del procesador de $\phi_1$ (segundo caso), hemos probado que $VC(\phi_1) < VC(\phi)$. De manera similar, sabemos que $VC(\phi) < VC(\phi_2)$. Como el orden definido para los relojes vectoriales es un orden parcial, en particular se cumple transitividad. Por lo tanto, $VC(\phi_1) < VC(\phi_2)$.
    \end{enumerate}
    
    \item Ahora supongamos que $VC(\phi_1) < VC(\phi_2)$ y veamos que $\phi_1 \rightarrow \phi_2$. \\
    Sea $\phi_1$ un evento de $p_i$ y $\phi_2$ un evento de $p_j$. Tenemos dos casos:
    \begin{enumerate}
        \item $i = j$ (mismo procesador). \\
        En particular tenemos que $VC(\phi_1)[i] \leq VC(\phi_2)[i]$ por hipótesis. Pero $VC(\phi_1)[i] \not= VC(\phi_2)[i]$ ya que la i-ésima entrada del vector $VC_i$ incrementa en uno con cada evento, por lo que necesariamente $VC(\phi_1)[i] < VC(\phi_2)[i]$. Como son del mismo procesador, esto sólo puede suceder si $\phi_1$ ocurre antes que $\phi_2$, es decir, $\phi_1 \rightarrow \phi_2$.
        \item $i \not= j$. \\
        En particular tenemos que $VC_i(\phi_1)[i] \leq VC_j(\phi_2)[i]$ por hipótesis. Además $VC_j(\phi_2)[i] \leq VC_i(\phi_1)[i]$ por la \textbf{Proposición 6.4} del libro de Hagit Attiya. En los enteros, esto implica que $VC_i(\phi_1)[i] = VC_j(\phi_2)[i]$. Pero la única forma en que $p_j$ pudo haber obtenido un valor para la i-ésima entrada de su vector que es al menos $VC_i(\phi_1)[i]$ es a través de una cadena de mensajes que originaron de $p_i$ en el evento $\phi_1$. Esto por definición significa que $\phi_1 \rightarrow \phi_2$.
    \end{enumerate}
    
\end{itemize}

\section{}

Supongamos que el corte consistente maximal no es único, es decir, que existen dos cortes consistentes maximales $s_1$ y $s_2$ tales que $s_1 \not= s_2$. \\
Por definición $s_1 \not< s_2$ (de lo contrario $s_1$ no sería maximal). Entonces tenemos dos casos:

\begin{itemize}

    \item $s_1 = s_2$. \\
    Esto es una contradicción porque supusimos que $s_1 \not= s_2$.
    
    \item $\exists k \in \{0,\ 1,\ ...,\ n - 1\}$ tal que $s_1[k] > s_2[k]$. \\
    Construimos otro corte $s$ de la siguiente manera: \\
    \[s[j] = \left\{\begin{array}{lr}
            s_1[j]\ \ \ si\ s_1[j] > s_2[j] \\
            s_2[j]\ \ \ en\ otro\ caso
            \end{array}\right\},\ \ \ para\ todo\ 0 \leq j \leq n - 1.\]
    Probemos que $s$ es un corte consistente. Tenemos que ver que el ($s[i] + 1$)-ésimo evento de $p_i$ no sucede antes que el $s[j]$-ésimo evento de $p_j$, para toda $i$ y $j$.
    Tenemos cuatro casos para el $i$-ésimo y $j$-ésimo elemento de $s$:
    \begin{enumerate}
        \item $s[i] = s_1[i]$ y $s[j] = s_1[j]$. \\
        Claramente el ($s[i] + 1$)-ésimo evento de $p_i$ no sucede antes que el $s[j]$-ésimo evento de $p_j$ ya que $s_1$ es consistente.
        \item $s[i] = s_2[i]$ y $s[j] = s_2[j]$. \\
        Claramente el ($s[i] + 1$)-ésimo evento de $p_i$ no sucede antes que el $s[j]$-ésimo evento de $p_j$ ya que $s_2$ es consistente.
        \item $s[i] = s_1[i]$ y $s[j] = s_2[j]$. \\
        Supongamos que el ($s_1[i] + 1$)-ésimo evento de $p_i$ sucede antes que el $s_2[j]$-ésimo evento de $p_j$. Por la construcción de $s$, $s_2[i] < s_1[i]$, de donde $s_2[i] + 1 < s_1[i] + 1$. Esto implica que el $(s_2[i] + 1)$-ésimo evento de $p_i$ sucede antes que el $s_2[j]-ésimo$ evento de $p_j$, lo cual es una contradicción porque $s_2$ es consistente.
        \item $s[i] = s_2[i]$ y $s[j] = s_1[j]$. \\
        Supongamos que el ($s_2[i] + 1$)-ésimo evento de $p_i$ sucede antes que el $s_1[j]$-ésimo evento de $p_j$. Por la construcción de $s$, $s_1[i] \leq s_2[i]$, de donde $s_1[i] + 1 \leq s_2[i] + 1$. Esto implica que el $(s_1[i] + 1)$-ésimo evento de $p_i$ sucede antes que el $s_1[j]-ésimo$ evento de $p_j$, lo cual es una contradicción porque $s_1$ es consistente.
    \end{enumerate}
    Por lo tanto, $s$ es un corte consistente. \\
    Como $s[j]$ se define como el máximo entre $s_1[j]$ y $s_2[j]$, entonces $s_1 \leq s$ y $s_2 \leq s$. Además $s_1 \not= s$ ya que si $s_1 = s$ esto implicaría que $s_2 < s_1$ ($s_1[j] > s_2[j]$ para toda $j$), lo cual no puede suceder porque $s_2$ es maximal. Luego, $s_2 \not= s$ porque sabemos que al menos $s[k] = s_1[k] > s_2[k]$. \\
    En resumen, suponiendo que el corte consistente maximal no es único hemos construido un corte consistente $s$ tal que $s_1 < s$ y $s_2 < s$, lo cual contradice la suposición de que $s_1$ y $s_2$ son maximales. \\
    Por lo tanto, el corte consistente maximal es único.
    
\end{itemize}

\section{}

Le pondré -1 si la línea pasa antes del primer evento (evento cero).

\begin{itemize}
    \item Primer corte: $s = [-1,\ -1,\ 0,\ 0]$. Es consistente porque no existen $i,j \in \{0,\ 1,\ 2,\ 3\}$ tales que $e_{i + 1, s[i] + 1}$ sucede antes que $e_{j, s[j]}$.
    \item Segundo corte: $s = [-1,\ 0,\ 0,\ ]$
\end{itemize}

(Ya me confundí porque se supone que un corte es un n-vector, pero en el diagrama las curvas pasan múltiples veces por cada proceso.) \\

Otra forma de determinar la consistencia es notando que si existe una flecha que comienza en el lado derecho de un corte y termina en el lado izquierdo, entonces el corte no es consistente. Dado esto, el primer y cuarto corte son consistentes, el segundo y tercero no lo son.

\section{}

En el instante en que le llega el mensaje de tomar un snapshot a $p_1$, los contadores de eventos son como sigue: \\
$num_1 = 1$ (envió un mensaje) \\
$num_2 = 0$ (aún no recibe el mensaje de $p_1$) \\
$num_3 = 1$ (envió un mensaje) \\
$num_4 = 0$ (aún no recibe el mensaje de $p_3$) \\
Entonces $p_1$ fija $ans_1 = num_1 = 1$ y envía un marcador a los demás, y las respuestas quedan así: \\
$ans_2 = 0$ \\
$ans_3 = 1$ \\
$ans_4 = 0$ \\
Luego $p_2$ envía marcadores a los demás pero no hacen nada porque $ans_i$ ya está fijo. Es decir, se envían 12 marcadores en total. La respuesta (el corte) queda como $[1,\ 0,\ 1,\ 0]$.

\end{document}
